% Part: first-order-logic
% Chapter: beyond
% Section: modal-logics

\documentclass[../../../include/open-logic-section]{subfiles}

\begin{document}

\olfileid{fol}{byd}{mod}

\olsection{模态逻辑}

考虑下面这个条件句的例子：

\begin{quote}
  如果 Jeremy 独自在那个房间里，那么他喝醉了，光着身子在椅子上跳舞。
\end{quote}

这是一个条件断言的例子，它可能实质为真但仍然具有误导性，因为它似乎暗示前件和结论之间存在着比单纯“前件为假或结论为真”更强的联系。也就是说，其措辞暗示该主张不仅在这个特定的世界中为真（在此它可能是平凡真的，因为 Jeremy 并非独自在房间），而且，更进一步，\emph{如果}前件为真，结论也\emph{会}为真。换言之，我们可以认为该断言意味着这个主张不仅在这个世界中为真，而且在任何“可能的”世界中都为真；或者说，它是\emph{必然}为真的，而不仅仅是在这个特定世界中为真。

模态逻辑正是为理解这种必然性而设计的。在普通命题逻辑中加入一个方框算子，就得到了模态命题逻辑；也就是说，如果 $!A$ 是一个公式，那么 $\Box !A$ 也是一个公式。直观上，$\Box !A$ 断言 $!A$ 是\emph{必然}为真的，或者说在任何可能世界中都为真。$\Diamond !A$ 通常被定义为 $\lnot \Box \lnot !A$ 的简写，可以读作断言 $!A$ 是\emph{可能}为真的。当然，模态也可以加入到谓词逻辑中。

克里普克结构可以用来为模态逻辑提供语义；事实上，克里普克 最初设计这种语义时考虑的就是模态逻辑。更一般地，我们不局限于偏序，而是有一个“可能世界”的集合 $P$，以及世界之间的二元“可达”关系 $\Atom{R}{x,y}$。直观上，$\Atom{R}{p,q}$ 断言世界 $q$ 与 $p$ 是相容的；也就是说，如果我们“处于”世界 $p$ 中，我们就必须考虑世界本可能像 $q$ 那样的可能性。

模态逻辑有时被称为一种“内涵”逻辑，以区别于“外延”逻辑。像经典逻辑这样的外延逻辑，其预期语义只涉及单一世界，即“实际”世界；而“内涵”逻辑的语义则依赖于更精细的本体论。除了用来刻画必然性，模态还可以用来刻画其他语言结构，根据应用领域重新解释 $\Box$ 和 $\Diamond$。例如：

\begin{enumerate}
\item 在可证性逻辑中，$\Box !A$ 读作“$!A$ 是可证的”，而 $\Diamond !A$ 读作“$!A$ 是一致的”。
\item 在认知逻辑中，人们可能将 $\Box !A$ 读作“我知道 $!A$”或“我相信 $!A$”。
\item 在时态逻辑中，可以将 $\Box !A$ 读作“$!A$ 总是真的”，而 $\Diamond !A$ 读作“$!A$ 有时是真的”。
\end{enumerate}

我们希望能够用处理模态的规则和公理来扩充逻辑。例如，系统 $\Log{S4}$ 由普通命题逻辑的公理和规则以及以下公理组成：
\begin{align*}
& \Box (!A \lif !B) \lif (\Box !A \lif \Box !B)\\
& \Box !A \lif !A \\
& \Box !A \lif \Box \Box !A
\intertext{以及一条规则：“由 $!A$ 推出 $\Box !A$。” $\Log{S5}$ 则增添以下公理：}
& \Diamond !A \lif \Box \Diamond !A
\end{align*}
这些公理的变体可能适用于不同的应用；例如，S5 通常被认为刻画了逻辑必然性的概念。而美妙之处在于，我们通常能通过以自然方式限制 克里普克 结构中的可达关系，来找到一个语义，使得推导系统相对于它是可靠且完备的。例如，$\Log{S4}$ 对应于可达关系是自反且传递的 克里普克 结构类。$\Log{S5}$ 对应于可达关系是{\em 普遍}的 克里普克 结构类，也就是说，每个世界都可以从其他任何一个世界到达；因此 $\Box !A$ 成立当且仅当 $!A$ 在每一个世界中都成立。

\end{document}